\documentclass[../../../include/open-logic-section]{subfiles}

\begin{document}

\olfileid[hi]{sth}{replacement}{ref}
\olsection{प्रतिस्थापन और परावर्तन}

\stagesinex{} द्वारा प्रतिस्थापन का औचित्य देने का हमारा अंतिम प्रयास एक
गहरे और सुंदर परिणाम से आरंभ होता है:\footnote{स्मरण रहे: सभी सूत्रों में
प्राचल हो सकते हैं (जब तक स्पष्टतः अन्यथा न कहा गया हो)।}
%In this, I will use overlining, such as $x_1, \ldots, x_n$, to
%abbreviate ``$x_1, \ldots, x_n$'': 
\begin{thm}[परावर्तन योजना]\ollabel{reflectionschema}
किसी भी सूत्र $\phi$ के लिए:
\[
\forall \alpha \exists \beta > \alpha (\forall x_1 \ldots, x_n \in
V_\beta)(\phi(x_1, \ldots, x_n) \liff \phi^{V_\beta}(x_1, \ldots, x_n))
\]
\end{thm}
\noindent 
\olref[sth][replacement][strength]{formularelativization} की तरह
$\phi^{V_\beta}$, $\phi$ के प्रत्येक परिमाणक को समुच्चय~$V_\beta$ तक सीमित
करने का परिणाम है। अतः सहजतः परावर्तन यह कहता है: यदि $\phi$ पूरे पदानुक्रम
में सत्य है, तो $\phi$ पदानुक्रम के मनमाने रूप से अनेक \emph{आरंभिक खंडों}
में सत्य है।

\citet{Montague1961} और \citet{Levy1960} ने दिखाया कि प्रतिस्थापन और
परावर्तन (के उपयुक्त सूत्रीकरण)~$\Z$ के सापेक्ष समतुल्य हैं, जिससे दोनों में
से किसी को जोड़ने पर~$\ZF$ मिलता है। (ये परिणाम हम
\olref[sth][replacement][refproofs]{sec} में सिद्ध करते हैं।) इस समतुल्यता
को देखते हुए कोई \stagesinex{} द्वारा परावर्तन और प्रतिस्थापन का औचित्य इस
प्रकार देने की आशा कर सकता है: \stagesinex{} को मानकर पदानुक्रम अत्यंत ऊँचा
होना चाहिए; वास्तव में इतना ऊँचा कि उसके विषय में कही जा सकने वाली कोई बात
उसकी ऊँचाई को परिबद्ध करने के लिए पर्याप्त न हो। इसे हम इस विचार के रूप में
समझ सकते हैं कि यदि कोई वाक्य~$\phi$ पूरे पदानुक्रम में सत्य है, तो वह
पदानुक्रम के मनमाने रूप से अनेक आरंभिक खंडों में सत्य है। और यही परावर्तन है।

फिर यह प्रतिस्थापन का आंतरिक औचित्य देने का सचमुच आशाजनक प्रयास लगता है।
किंतु यहाँ इसके विषय में कहने के लिए बहुत अधिक है। अब आपको स्वयं निर्णय करना
होगा कि यह सफल है या नहीं।\footnote{फिर भी आप
\citet[95--100]{Incurvati2020} पढ़कर आगे बढ़ना चाह सकते हैं।}

\end{document}
